\clearpage
\section{Scope, Edition, and Qualifications}
\textbf{Governing edition and jurisdiction.} Entries in the comprehensive adult reference dictionary describe the Roman Rite as governed by the \textit{Missale Romanum}, \textit{editio typica tertia} (2002; \textit{reimpressio emendata}, 2008), and its \textit{Institutio Generalis Missalis Romani} (GIRM), together with the approved English edition for the United States of America, \textit{The Roman Missal, Third Edition} (effective First Sunday of Advent, 27 November 2011), the adaptations approved by the United States Conference of Catholic Bishops (USCCB), the 2021 emendations incorporated in the official USCCB GIRM (Prot.~n.~228/20; Prot.~n.~394/21, 22 October 2021), and the 2021 amendment of canon 230 \S1 (\textit{Spiritus Domini}, 10 January 2021). Norms are stated as of 2026-11-25.

\textbf{Coverage and boundaries.} Within this reference corpus, 40 principal sanctuary furnishings, sacred vessels, linens, liturgical books, service items, and vestments of Mass and closely joined rites are documented alongside 3 historically superseded postconciliar provisions. Episcopal and pontifical insignia (miter, crosier, pectoral cross, ring, zucchetto, gremiale, bugia), full Paschal Triduum apparatus (Paschal candle, Easter fire brazier, Good Friday adoration cross veil), funeral pall and catafalque, and specialized rites of dedication and ordination remain outside the present scope and are reserved for future expansion rather than treated by conjecture.

\textbf{Separation of universal law, territorial adaptation, and custom.} Universal provisions of the Latin typical edition (including choices entrusted to the Diocesan Bishop, such as tabernacle placement at GIRM 315), approved United States adaptations (solid wood fixed altars at GIRM 301, festive under-cloths at GIRM 304, sacred-vessel materials at GIRM 329, lay ministerial attire at GIRM 339, Communion under both kinds at GIRM 283 and the USCCB Norms, and funeral and solemn vestment colors at GIRM 346e and 346h), diocesan or celebrant options, and documented parish custom are distinguished throughout. Forms belonging to the 1962 Roman Missal appear only in explicit comparison fields and keyed source notes, never as postconciliar norms.

\textbf{Rights, visual provenance, and reliance boundary.} All ten monochrome graphite plates are original project-created illustrations normalized to 8-bit grayscale on pure white ground; no third-party scans, photographs, or copyrighted illustrations are reproduced. English descriptions paraphrase official liturgical law in original prose without reproducing protected ICEL liturgical prayer texts or rubrical formulations verbatim. English Scripture quotations follow the public-domain Douay--Rheims (Challoner) text. The approved Roman Missal and diocesan directives remain the governing liturgical law; study and preparation notes here carry no official, liturgical, or magisterial authority.

\phantomsection
\section*{References}
\addcontentsline{toc}{section}{References}
\subsection*{Governing Postconciliar Roman Liturgical Books and Law}
\begin{itemize}
\item Missale Romanum, ex decreto Sacrosancti Oecumenici Concilii Vaticani II instauratum, auctoritate Pauli PP. VI promulgatum, Ioannis Pauli PP. II cura recognitum, editio typica tertia (Typis Vaticanis, 2002); reimpressio emendata (Typis Vaticanis / Libreria Editrice Vaticana, 2008; announced in Notitiae 44 [2008], nos. 503--504, pp. 367--388), including Institutio Generalis Missalis Romani (IGMR).
\item The Roman Missal, Third Edition, for Use in the Dioceses of the United States of America (Washington, DC: United States Conference of Catholic Bishops, 2011; CDWDS recognitio 26 March 2010 and 24 July 2010; effective 27 November 2011), including the General Instruction of the Roman Missal with U.S. adaptations and 2021 emendations (Prot. n. 228/20; Prot. n. 394/21, 22 October 2021).
\item United States Conference of Catholic Bishops, Norms for the Distribution and Reception of Holy Communion under Both Kinds in the Dioceses of the United States of America (approved 15 June 2001; CDWDS recognitio 22 March 2002, Prot. 1383/01/L; revised 2011; reflecting October 2006 CDWDS decision on vessel purification).
\item United States Conference of Catholic Bishops, Excerpts from the Roman Missal (Book of the Chair; CDWDS confirmed; Washington, DC: USCCB, 2018).
\item Congregation for Divine Worship and the Discipline of the Sacraments, Response on GIRM no. 299 (25 September 2000, Prot. N. 2036/00/L; Notitiae 36 [2000], 171--173); Instruction Redemptionis Sacramentum (25 March 2004; AAS 96 [2004], 549--601); Decree implementing the provisions of canon 838 of the Code of Canon Law, Postquam Summus Pontifex (22 October 2021).
\item Ordo lectionum Missae, editio typica altera (Libreria Editrice Vaticana, 21 January 1981), Praenotanda nos. 1--125.
\item Caeremoniale Episcoporum, ex decreto Sacrosancti Oecumenici Concilii Vaticani II instauratum, editio typica (Typis Polyglottis Vaticanis, 14 September 1984).
\item Codex Iuris Canonici auctoritate Ioannis Pauli PP. II promulgatus (25 January 1983; AAS 75/II [1983]), cann. 230 (as amended by Pope Francis, Motu Proprio Spiritus Domini, 10 January 2021), 847, 858, 910, 938--940, 1003, 1235--1239.
\item Catechism of the Catholic Church, 2nd ed. (Vatican City: Libreria Editrice Vaticana, 1997; Washington, DC: USCCB, 2019).
\item Sacred Congregation for Divine Worship, Ordo Baptismi Parvulorum (Rite of Baptism for Children; editio typica, 1969; editio typica altera, 1973); Ordo benedicendi oleum catechumenorum et infirmorum et conficiendi chrisma (3 December 1970); De sacra communione et de cultu mysterii eucharistici extra Missam (21 June 1973).
\end{itemize}
\subsection*{Earlier Postconciliar Legislative Milestones (1963--1975)}
\begin{itemize}
\item Second Vatican Council, Constitution on the Sacred Liturgy Sacrosanctum Concilium (4 December 1963; AAS 56 [1964], 97--138).
\item Sacred Congregation of Rites and Consilium, Instruction Inter Oecumenici (26 September 1964; AAS 56 [1964], 877--900), Instruction Tres abhinc annos (4 May 1967; AAS 59 [1967], 442--448), and Instruction Eucharisticum Mysterium (25 May 1967; AAS 59 [1967], 539--573).
\item Pope Paul VI, Apostolic Constitution Missale Romanum (3 April 1969; AAS 61 [1969], 217--222); Motu Proprio Mysterii Paschalis (14 February 1969; AAS 61 [1969], 222--226); Motu Proprio Ministeria quaedam (15 August 1972; AAS 64 [1972], 529--534).
\item Sacred Congregation for Divine Worship, Missale Romanum, editio typica prima (26 March 1970, Prot. 166/70) and editio typica altera (27 March 1975); Sacred Congregation for the Discipline of the Sacraments, Instruction Immensae caritatis (29 January 1973; AAS 65 [1973], 264--271).
\end{itemize}
\subsection*{1962 Comparison Witnesses, Scripture, and Liturgical Tradition}
\begin{itemize}
\item Missale Romanum, editio typica (Typis Polyglottis Vaticanis, 1962), including Rubricae Generales Missalis Romani (1960/1962) and Ritus Servandus in Celebratione Missae; together with Codex Iuris Canonici (1917), Rituale Romanum (1925/1952), and Pontificale Romanum (1859).
\item Adrian Fortescue, The Ceremonies of the Roman Rite Described (London: Burns \& Oates, 1917; 7th rev. ed., 1943).
\item The Holy Bible, Douay--Rheims Version (Challoner revision).
\item St. Cyril of Jerusalem, Mystagogical Catecheses (II.4); St. Thomas Aquinas, Summa Theologiae, Tertia Pars, q. 83, aa. 3--5 (on sacred vessels, linens, vestments, and altar rites); Amalarius of Metz, Liber Officialis (De ecclesiasticis officiis, PL 105); traditional Roman vesting prayers of the Missale Romanum.
\end{itemize}
